\section*{Updated overlap interface}

Let
\[
 q(y)=y\log_2(1/y)+\frac{2y}{\ln2},\qquad
 \Phi(f)=\sum_i q(\widehat f(\{i\})^2).
\]
For a distinguished physical coordinate $x$, define
\[
 \Lambda_x(f)=q(\widehat f(\emptyset)^2+\widehat f(\{x\})^2)
 +\sum_{i\ne x}q(\widehat f(\{i\})^2+
                         \widehat f(\{i,x\})^2).
\]

The following arrows are proved:
\[
 \begin{array}{c}
 f(x,u),g(x,v),\ u\cap v=\emptyset
 \\
 \Downarrow\quad\text{one shared physical coordinate}
 \\
 \Phi(fg)\le\Lambda_x(f)+\Lambda_x(g),
 \qquad
 \mathbb E_{x}\Phi(f\mid x)\le\Lambda_x(f).
 \end{array}
\]
The first arrow is stable under arbitrary private tensor factors because the
private blocks $u,v$ are unrestricted Boolean systems; no rank, bias,
dimension, symmetry, or rarity hypothesis is used.

The B05/A025 counterfamily proves that replacing $\Lambda_x(f)$ by
$\Phi(f)+C\operatorname{Inf}_x(f)$ is impossible at fixed $C$. Thus the block
term is genuinely stronger information, not a restatement of the failed
scalar payment.

The next arrow is open:
\[
 \left.
 \begin{array}{c}
 \text{adaptive hierarchy with reused coordinates}\cr
 \text{each }\Lambda_x\text{ owned at first physical revelation}
 \end{array}
 \right\}
 \stackrel{\rm OPEN}{\Longrightarrow}
 \sum_{\rm nodes}\text{signed overlap charge}
 \le A I(f)+bU(f),\qquad b<1.
\]
If this open arrow holds with absolute $A$ and $b<1$, the existing restriction
induction gives
\[
 U(f)\le\frac{A}{1-b}I(f),\qquad
 H(f)\le\left(2+\frac{2A}{1-b}\right)I(f).
\]

The open arrow must retain the exact mismatch/KL term and prevent one Fourier
block from being charged by multiple ancestors. A fixed multiplier on each
shared pivotal edge is already refuted.
